\documentclass[11pt,reqno]{amsart}
\usepackage[T1]{fontenc}
\usepackage{lmodern}
\usepackage{amssymb,array,booktabs}
\usepackage[margin=1.15in]{geometry}
\setlength{\emergencystretch}{1.3em}
\usepackage[hidelinks]{hyperref}
\hypersetup{pdftitle={Eleven-digit Kaprekar arrival laws modulo twelve},pdfauthor={Rustem Galyavin}}
\newtheorem{theorem}{Theorem}
\newtheorem{proposition}[theorem]{Proposition}
\newtheorem{lemma}[theorem]{Lemma}
\theoremstyle{definition}
\newtheorem*{definition}{Definition}
\newcommand{\Prob}{\mathbb P}
\newcommand{\Z}{\mathbb Z}
\newcommand{\R}{\mathbb R}
\newcommand{\calY}{\mathcal Y}
\newcommand{\vct}{\mathbf v}
\newcommand{\cen}{c_*}
\newcommand{\one}{\mathbf 1}
\newcommand{\T}{T_B}
\newcommand{\ZB}{Z_B}
\newcommand{\Id}{\mathrm{Id}}
\newcommand{\Zhat}{\widehat{\Z}}

\title[Eleven-digit Kaprekar arrival laws]{Eleven-digit Kaprekar arrival laws modulo twelve}
\author{Rustem Galyavin}
\thanks{\href{mailto:rustem@galyavin.com}{rustem@galyavin.com}. ORCID \href{https://orcid.org/0009-0006-6875-9170}{0009-0006-6875-9170}.}
\date{20 November 2026}
\subjclass[2020]{11A63, 37P99, 60B10}
\keywords{Kaprekar map, arrival time, weak convergence, digit differences, lattice index}
\begin{document}
\raggedbottom
\begin{abstract}
Choose an eleven-digit word uniformly in base \(B\), allowing leading zeros and
excluding constant words. Let \(T_B\) be the number of Kaprekar subtractions
to the first cyclic word. Couple all bases through eleven uniform real digits.
On the event that the leading orbit of their ordered differences enters a
five-phase open chamber, an event of probability at least
\(92.9921494990\%\), the normalized initial digit differences, \(T_B/B\), and
the normalized differences of the first cyclic word have a joint weak limit
along every base class modulo twelve. For a positive integer \(M\), these
measures converge along every class modulo \(M\) if and only if twelve
divides \(M\). A clock kernel describes the limit by charging each repeated
affine return its limiting duration and giving each admissible exit residue
a probability equal to the reciprocal of its preimage index in the inherited
input lattice. Although some rank-one arrivals repeat only with period
sixty-six in \(B\), the extra factor eleven disappears from the joint law.
The proof uses exact polyhedral and lattice computations.
\end{abstract}
\maketitle

\section{The entrant law}

A Kaprekar subtraction takes the difference between the decreasing and
increasing rearrangements of the digits. Keep all eleven slots, and count time
zero if the original word is already cyclic. For sorted digits
\(d_1\ge\cdots\ge d_{11}\), the difference vector has entries
\(q_i=d_i-d_{12-i}\), \(1\le i\le5\). Write \(q_B\) for this vector at time
zero, \(\T\) for first cyclic-word time, and \(\ZB\) for the difference vector
of that word divided by \(B\). The observable thus records the first cyclic
word through its five differences.

Couple all bases by drawing eleven independent uniform digits from \([0,1)\)
and flooring their products with \(B\). Write \(y\) for their real ordered
differences and put
\[
 \calY=\{1\ge y_1\ge\cdots\ge y_5\ge0\},\qquad
 w(y)=11!\prod_{i=0}^5g_i(y),
\]
\[
 g(y)=(1-y_1,y_1-y_2,\ldots,y_4-y_5,y_5).
\]
We write \(m(A)=\int_A w(y)\,dy\) for the real-input probability. Let \(K\)
sort \(1,y_i,1-y_i\) and take the five symmetric differences. To describe the
open chamber \(\Pi\), put
\[
 L(y)=(y_1,y_1-y_5,y_2-y_5,y_2-y_4,y_3-y_4).
\]
At each of \(y,Ly,\ldots,L^4y\), require positive gaps and the strict order
\[
 1>y_1>1-y_5>y_2>1-y_4>\cdots>y_5>1-y_1 .
\]
On this chamber \(K=L\), and on gaps \(L\) permutes the last five entries by
\((g_1,g_2,g_3,g_4,g_5)\mapsto(g_5,g_1,g_4,g_2,g_3)\) with \(g_0\) fixed.
In particular \(K^5y=y\) there. Set
\[
 D=\{y:K^j y\in\Pi\text{ for some }j\ge0\},\qquad
 \vct=(5,4,3,2,1),\qquad \cen=\vct/6.
\]
The event \(\{y\in D\}\) is defined by the coupled real digits, not by the
base-\(B\) word. The terms used in the next definition are fixed in
Section~\ref{sec:routes}.

\begin{definition}[Clock kernel]
A compatible residue system \(\beta\in\Zhat\) gives one residue of \(B\)
modulo every positive integer, consistent under reduction. For \(y\in D\)
off the null walls of Section~\ref{sec:passage}, let \(j(y)\) be the least
\(j\) with \(K^jy\in\Pi\), and let \(\omega\) be Haar-uniform on
\(\Zhat^{\,5}\), independent of \(y\); it plays the role of the residue of
the original integer input. From the entrance chart of \(y\), follow the
physical route of Section~\ref{sec:routes}. At a translating return of
\(\ell\) updates whose winning guard \(G\) has decrement \(s>0\), read the
exit residue \(G\bmod s\) from \((\beta,\omega)\) and add \(\ell\alpha/s\)
to the clock, where \(\alpha=\lim G/B\); bounded primitive switches add
nothing. Let \(\tau_\eta\) be the clock at the first visited section of
leading radius less than \(\eta\). Then \(\tau_\eta\) is nondecreasing as
\(\eta\downarrow0\), and \(\tau_{\eta'}\le\tau_\eta+(16/3)\eta\) for
\(\eta'<\eta\) (Section~\ref{sec:kernel}). Put
\(\tau=\lim_{\eta\downarrow0}\tau_\eta\) and let \(Q_\beta(y,\cdot)\) be
the law of \(\tau\). For \(r\in\Z/12\Z\), let \(Q_r=Q_\beta\) with \(\beta\)
the integer representative \(0\le r<12\).
\end{definition}

\begin{theorem}\label{th:main}
For each \(r\in\Z/12\Z\), as \(B\to\infty\) through \(B\equiv r\pmod{12}\),
\[
 \Prob\bigl(y\in D,\ (q_B/B,\T/B,\ZB)\in\cdot
       \,\bigm|\,\text{nonconstant word}\bigr)
 \ \Longrightarrow\
 \one_D(y)w(y)\,dy\,Q_r(y,dt)\,\delta_{\cen}(dz)
\]
weakly as finite measures, with \(Q_r\) the clock kernel defined above.
Consequently, for a positive integer \(M\), these measures converge along
every residue class modulo \(M\) if and only if \(12\mid M\).
\end{theorem}

\begin{proposition}\label{prop:mass}
Let \(E_3\) require that at time three both end gaps are at least \(1/8\), the
sum of the three smallest gaps is at most \(3/8\), and the largest gap is at
least \(1/4\). Then
\[
 \begin{split}
 m(E_3)&\le 0.65165187185\%,\\
 m(D)&\ge 1-\frac{5124467}{80621568}-m(E_3)\ge 92.9921494990\%.
 \end{split}
\]
The probabilities are relative to the original eleven independent digits.
\end{proposition}

\cite[Theorem~1]{G6} proves tightness of \(\T/B\) and convergence in probability
of \(\ZB\) on the whole input population. \cite[Theorem~4]{G7} proves two
different laws on a box of original mass about \(8.5\cdot10^{-14}\), and that any
common modulus for the full joint law must be divisible by twelve; whether
that law converges along the classes modulo twelve is left open there.
Theorem~\ref{th:main} supplies a system of joint limits on an entrance event
of mass at least \(92.99\%\), with smallest common modulus exactly twelve.

The divisibility in \cite{G7} does not by itself give minimality here,
because the measures of Theorem~\ref{th:main} test only inputs in \(D\). The
midpoint of the box \(V\) of \cite[Theorem~4]{G7} is a strict interior point
of \(\Pi\), and the comparison of classes \(5\) and \(11\) there gives the
factor four; we recompute it with the kernels \(Q_r\). The six laws of
\cite[Theorem~2]{G7}, which force divisibility by six there, live on a
rectangle that is not contained in \(\Pi\) at time zero. Rather than transport them
through an entry prefix, we compare classes \(9\) and \(5\) on a small cube
inside \(\Pi\) (Lemma~\ref{lem:local}).

\emph{Why the base residue enters.} After entry into \(\Pi\) the integer
orbit spends almost all of its time in translating returns: a fixed word of
\(\ell\) subtractions moves the state by a fixed integer vector and is
repeated while all its guards hold. A guard \(G\) with decrement \(s\) fails
after \(1+\lfloor G/s\rfloor\) copies. Since \(G=\alpha B+O(1)\) is an
integral affine function of \(B\) and of the input, the return contributes
\(\ell\alpha/s\) to \(T_B/B\), while the exit state depends on \(G\bmod s\),
that is, on the residues of \(B\) and of the input. Averaging over inputs
turns each admissible residue into a weight, the reciprocal of its lattice
index. The proof has four steps. (i) After at most six macroscopic returns
the direction space has rank one (Lemma~\ref{lem:ingress}). (ii) The outer rank-one graph is finite and acyclic, and each inner arrow
contracts the leading radius by a factor at most \(251/256\), so every
positive collar is reached. At a fixed section in the \(\eta\)-collar,
capture still ahead takes at most \((16/3)\eta B+O(1)\) further updates;
if capture has occurred, the section follows it by \(O_{\rm record}(1)\)
updates (Section~\ref{sec:passage}).
(iii) On paired states the centred dynamics does not depend on \(B\)
\cite[Theorem~4]{G4}, so one modulus governs every depth
(Section~\ref{sec:onemod}). (iv) The periods that occur are \(6\), \(12\)
and \(66\), and the factor eleven cancels because the period-\(66\) returns
end at the centre with a leading duration that does not depend on the
phase (Section~\ref{sec:eleven}). Steps (i), (ii) and (iv) rest on finite
exact certificates (Appendix~\ref{ap:certificates} and
Section~\ref{sec:verification}).

Devlin and Zeng \cite[Theorem~1.1]{DZ} determine maximum distances to the
four-digit fixed point. Four-digit cycle classifications include
\cite{G1,COST}, and Kay and Downes-Ward \cite{KDE} classify families of
fixed points and cycles in even bases through arithmetic modulo \(B-1\).
Thakur \cite[Section~2.7]{Thakur} gives the borrowed digit representation
used below. Dahl \cite{Dahl} studies gap-based transition statistics at
fixed base ten. Growing-base arrival laws for four, five, seven and nine
slots appear in \cite{G2,G3,G5,G4}. We use the odd-length majorization
theorem of \cite{G4}, the leading and finite estimates of \cite{G6}, and
the local clock laws of \cite{G7}.

\section{Exact finite routes and their weights}\label{sec:routes}

\emph{Terminology.} \emph{Physical} refers to the exact integer orbit of
the difference update \(F_B\) at a finite base; \emph{leading} refers to its
limit after division by \(B\), governed by \(K\), the continuous extension
of the large-base limit of \(F_B/B\) \cite[Section~2]{G6}. A \emph{cyclic
word} lies on a cycle of the Kaprekar map, and \emph{capture} is arrival at
the first cyclic word, at time \(T_B\). A \emph{root} is a class of
\(B\bmod6\); it fixes the entrance lattice \(X_i\equiv iB\pmod6\) of
Section~\ref{sec:entry}. \emph{Ingress} is the entry into \(\Pi\). A
\emph{primitive update} is one subtraction on a fixed sorting branch. A
translating return (Section~\ref{sec:returns}) is \emph{macroscopic} when
its number of copies grows linearly in \(B\), and a \emph{primitive switch}
is a primitive update outside the complete copies of such returns. The
\emph{rank} is the dimension of the real direction space \(V\) of the
current affine sheet. A \emph{compatible base progression} is an arithmetic
progression of bases on which all congruences of a record are consistent.
A \emph{section} is a rank-one sheet with an interval of leading positions
(Section~\ref{sec:collar}); it is \emph{inner} if it lies strictly inside
\(R(X)/B<1/7\) and \emph{outer} otherwise, and a \emph{collar} is a region
\(R(X)/B<\eta\). The \emph{permanent region} is
\(\{\min G\ge\lceil(B+1)/7\rceil+64\}\), in which a paired state stays
paired for \(B\ge3534\) \cite[Section~6]{G6}.

\subsection{Borrowing and entry}\label{sec:entry}
The twelve spacings of eleven uniform order statistics have constant
density \(11!\). The six gaps are sums of pairs of spacings. Integrating
over the six splits gives \(w\). For integers put \(G=B g(q/B)\).
The number of sorted words of differences \(q\) is
\[
 (B-q_1)\prod_{i=1}^5(q_i-q_{i+1}+1),\qquad q_6=0.
\]
A tie fixes one part of a split and gives at most \(O(B^5)\) choices.
Consequently the number of ordered words is \(B^6w(q/B)+O(B^5)\),
uniformly on strict input cells. The sample space has \(B^{11}-B\)
nonconstant words. On the common real input the conditioning changes
probabilities by \(O(B^{-10})\), and \(\|q_B/B-y\|_\infty<1/B\).

Let \(F_B\) denote the exact difference update. At positive height
\(h=\#\{i:q_i>0\}\), the digit to the right is \(B-q_1\); the other
negative columns give \(B-q_i-1\) for \(2\le i\le h\). The \(11-2h\)
middle columns give \(B-1\), and the positive columns give \(q_i\) for
\(i<h\) and \(q_h-1\) at the innermost positive column. Sort this full
list to compute \(F_B\). For \(h=5\) the five borrowed levels, following
the singleton \(B-1\), are
\[
 \bigl(q_i-\one_{i=5},\ B-q_{6-i}-\one_{i<5}\bigr),
       \qquad i=1,\ldots,5.
\]
We say a state is paired when each level lies weakly above the next.
Full sorting is always used to define the numerical update. We fix a
deterministic order for equal digits. The weak paired comparison is
recorded separately from the tie-resolved sorting comparisons.

Each strict leading branch is unimodular. Indeed for
\(r_1>\cdots>r_5>0\) the sorted magnitudes of \(2y_i-1\), the output is
\[
 K(y)=\tfrac12(1+r_1,r_1+r_2,r_2+r_3,r_3+r_4,r_4+r_5).
\]
The derivative of the magnitude ordering is twice a signed permutation,
and the adjacent-sum matrix divided by two cancels the factor two in
each dimension. On a compact subset of any strict finite entry prefix,
the physical branches have the same derivative and bounded offsets.
Thus the affine entry map sends the unrestricted original integer
lattice onto the full entrance lattice before imposing the transported
cell. In centered coordinates \(X=6q-B\vct\), this lattice is
\(X_i\equiv iB\pmod6\). Proper leading walls are null, including their
finite leading preimages. Each entry chart retains its own original
position and density.

Put
\[
 R(X)=-\min(-X_1,X_1-X_2,\ldots,X_4-X_5,X_5),\qquad
 e(y)=\sum_{i=0}^5(g_i(y)-1/6)^2 .
\]
Then \(\min G=(B-R(X))/6\). The closure of \(\Pi\) has 77 vertices,
and its only maximizer of \(e\) is the boundary gap vector
\((0,1/5,1/5,1/5,1/5,1/5)\), of energy \(1/30\)
(Lemma~\ref{lem:ingress}). On compact subsets of strict entry charts, \cite[Section~5]{G6} gives
\(q_5>0\) (height five) until entry into the permanent paired region, and
permanent positive pairing gives \(q_5>0\) thereafter. This supplies the
hypothesis for all five borrowed levels.

\subsection{Repeated returns and residue splits}\label{sec:returns}
A chart records an open polyhedral set of original normalized inputs,
a compatible affine lattice \(q_0\in p_B+\Lambda\) of original integer
inputs, and the state \(X=Mq_0+B\mathbf b+c\). A \emph{guard} is an integer affine
inequality \(G_i(X,B)\ge0\) specifying a primitive branch, pulled back to
the start of the proposed return. An integer strict comparison uses a
unit offset. A certified word of \(\ell\) primitive updates fixes the
leading chart. Its exact affine map translates this sheet by a
nonzero integer vector \(d\), and its linear part \(A\) satisfies
\(Ad=d\); we call it a \emph{translating return}. Every complete copy
then advances by \(d\), provided its
guards hold. The guard with state row \(a_i\) changes by \(a_i d\) per
copy. The number of complete copies is
\[
 1+\left\lfloor\min_{s_i>0}\frac{G_i(X,B)}{s_i}\right\rfloor,
 \qquad s_i=-a_i d .
\]
After this translation the next primitive phases are read at the exact
state, including any valid partial copy before the first failed guard.
Full sorting permits a return word to contain unpaired steps. When
two time forms tie identically on the chart, their difference
annihilates its variations and its leading base coefficient cancels.
Compare their physical intercepts. A nonidentical leading tie is a
proper original wall.

On a cell where \(G\) wins with decrement \(s>0\), let
\(\alpha=\lim G/B\) and \(j=G\bmod s\). The allowed residues form one coset
of the image of \(a_G M\Lambda\) in \(\Z/s\Z\), where \(a_G\) is the
state row of \(G\). For \(\nu=\gcd(s,a_G M\Lambda)\), each of the \(s/\nu\)
possible residues has preimage of relative index \(s/\nu\). The exit
is obtained after \(1+(G-j)/s\) copies. A choice of residue restricts the
original input lattice to that preimage, sometimes called a \emph{lane}.
Subsequent congruences are imposed on this preimage, even when its
image in the state space loses rank. Their relative probabilities
are the reciprocals of the corresponding indices. The scaled duration
of the return is \(\ell\alpha/s\). Each bounded sequence of primitive
switches between returns has zero limiting duration.

For a finite route, clear denominators in all affine charts and
congruences. On each compatible base progression a leaf is an
original cell \(P\) and an original affine lattice of index \(I\).
For any bounded continuous test of its stopped marks, weighted Riemann
sums give the contribution
\begin{equation}\label{eq:transfer}
 I^{-1}\int_P w(y) f(y,S(y),z(y))\,dy ,
\end{equation}
where \(S,z\) are the limiting elapsed time and the limiting marker
being carried. A fundamental parallelepiped scaled by \(1/B\) has
volume \(I/B^5\), and the point weights are \(B^{-5}w+O(B^{-6})\).
One can trim collars of proper walls first. The integer branches and
affine formulas then hold uniformly on the remaining compact sets.
Original cells stay marked when their images overlap or coincide.

\section{Rank passage and approximation of the first cyclic time}\label{sec:passage}

\subsection{Access to every collar}\label{sec:collar}
An exit with state guard row \(a_G\) has derivative
\(H=\Id+d a_G/s\). Since \(a_G d=-s\), its kernel is exactly \(\R d\).
A residue split preserves the real incoming direction space \(V\), so
the rank drops by one precisely when \(d\in V\). The six-root ingress
calculation (Lemma~\ref{lem:ingress}) gives the rank histories
\[
 54321,\qquad 543221,\qquad 5432221,
\]
through first rank one, where the digits list ranks after successive
macroscopic exits, including ingress. The total number of primitive
switches through ingress to rank one is at most 230 on even roots
and 246 on odd roots. Rank-two returns can preserve rank, and the
calculation includes unpaired primitive phases.

At rank one write the sheet as
\[
 X=p+kT+n\ell_0 u,\qquad B=b+kP,
\]
with primitive integer tangent \(u\). Choose an integral row \(c_0\)
with \(c_0u=1\), and put \(N=\Id-uc_0\). The exact normal equation is
\(NX=(NT/P)B+Np-bNT/P\). Since \(N(6\Z^5)=6(\Z^5\cap\ker c_0)\),
the least compatible normal base shift preserving the physical phase
is \(6\operatorname{lcm}_i\operatorname{den}((NT/P)_i)\).
For coverage one may enlarge the tangent to the full physical
tangent \(6u\Z\) on this sheet. Equation~\eqref{eq:transfer} still
uses the inherited lattice.

A rank-one section is a normal sheet with an interval of leading
positions. Call \(R(X)/B<\eta\) a terminal collar, with
\(0<\eta<1/7\). Sections strictly inside \(R(X)/B<1/7\) are permanent
for sufficiently large bases on compact subsets, because then
\(\min G>\lceil(B+1)/7\rceil+64\). The exact rank-one graphs of
Lemma~\ref{lem:lines} are finite, their outer subgraphs are acyclic,
and an inner successor stays inner or reaches the centre. The leading
lines are homogeneous about \(X/B=0\); the interior part of a line is
enlarged to its full open half-ray in \(R(X)/B<1/7\). Periodicity of
its leading seed transfers by positive dilation in that paired region,
and the guards are verified on the enlarged interval. Outer intervals
retain their actual locations.

By \cite[Section~6 and Appendix~E]{G6}, until the first \(R\le120\)
a lower hit by at least six units occurs within 32 updates, the
intermediate overshoot is at most 384, and after that hit \(R\le504\).
Consequently in the permanent region,
\[
 R_t\le\max(504,R_0-6\lfloor t/32\rfloor+384).
\]
An inner arrow of positive scaled duration \(\Delta\) thus satisfies
\(r'\le\max(0,r-3\Delta/16)\) for leading radii. Its certified exact
minimum ratio \(\Delta/r\) is \(5/48\), giving
\(r'\le(251/256)r\). Finite outer passage and inner contraction give
access to every positive collar after finitely many returns. The
exceptional walls over all finite records form a countable union of
original null sets. A proper wall in a projected chart pulls back
to a proper wall because the projection has full rank in its
retained coordinates.

\subsection{The section stop may follow capture}
Let \(S_B\) be the physical time of a fixed certified section strictly
inside a collar of radius \(\eta\). For large bases on a compact
margin, the remaining number of updates until capture, if capture
is ahead, is at most
\begin{equation}\label{eq:forward}
 32\max(0,\lceil(R-120)/6\rceil)+\binom{509}{5}+1
 \ \le\ (16/3)\eta B+O(1).
\end{equation}
Indeed, after the first hit of \(R\le120\), the orbit stays in
\(R\le504\). In this envelope the six centered gap offsets, after adding
504, are nonnegative integers summing to \(6\cdot504\). Their common
residue modulo six gives at most \(\binom{509}{5}\) possible states.
The extra update passes from a cyclic difference vector to a cyclic
word, since the next word is a function of that vector.

If capture precedes \(S_B\), the eventual cycle lies in that
envelope because the orbit has a later permanent section. Every
same-phase state on the cycle has \(\|X\|_\infty\le2520\). Consecutive
same-phase states in a certified block are separated by its nonzero
integral displacement \(d\). In a coordinate with \(d_i\ne0\) there
are at most 5041 consecutive such points after capture. Complete
blocks, bounded switches and the incomplete last block therefore
give only \(O_{\text{record}}(1)\) possible updates after capture.
This proves the two sides of the comparison
\begin{equation}\label{eq:two-sided}
 S_B-O_{\text{record}}(1)\le T_B\le S_B+(16/3)\eta B+O(1).
\end{equation}
The argument permits the first physical visit to the collar to
occur within an earlier repeated word.

\subsection{The clock kernel and weak convergence}\label{sec:kernel}
We check that the clock kernel of Section~1 is well defined. Fix a
compatible system \(\beta\), a real input in \(D\) off the null walls,
and its first strict leading entry prefix. Use the chronological
physical guards and draw every allowed remainder with the conditional
index rule of Section~\ref{sec:returns}. This defines consistent finite
cylinders because the children partition the parent coset. Equivalently
take a single Haar-uniform element \(\omega\) of \(\Zhat^{\,5}\),
independent of \(y\), for the original residue and read its successive
finite quotients. Declare a section with leading state zero terminal.
For \(\eta\downarrow0\), stop at the first visited section with leading
radius less than \(\eta\). The chronological clock forms at these stops
are nondecreasing. When a deeper stop follows one in the
\(\eta\)-collar, lift the two records on the same original chart and use
\eqref{eq:two-sided} at both stops. On dividing by \(B\) and taking the
compatible limit, the additional elapsed time is at most
\((16/3)\eta\). The limit is therefore finite and measurable. Its law
for each such \(y\) is the kernel \(Q_\beta(y,dt)\) of Section~1. This
definition continues to arbitrarily small collars, with no fixed cap on
entry time or number of returns.

Take a sequence of bases that is eventually in the class
\(\beta\bmod m\) for every fixed \(m\). Fix a bounded Lipschitz test and
let \(D_J=\bigcup_{j=0}^J K^{-j}(\Pi)\).
The finite entrance prefix costs at most \(J/B\). Under the common
real-input coupling, the omitted mass of \(D\setminus D_J\) is at most
\(m(D\setminus D_J)/(1-B^{-10})\), uniformly over bases. Fix a collar
radius away from atoms of the section radii and take finitely
many disjoint first-stop cylinders on
\(D_J\) that omit mass at most \(\varepsilon\) more. These can be trimmed
to compact strict subcells at arbitrarily small cost. Their finite
physical and limiting contributions agree asymptotically by
\eqref{eq:transfer}. Comparison with their total original mass also
bounds the limsup mass omitted in the finite-base measures. This
argument avoids any requirement on the boundary of \(D\).

On the retained cylinders the scaled clock error is at most
\((16/3)\eta+o(1)\), and \(\ZB\) converges to \(\cen\) by the cycle
envelope, or by \cite[Theorem~1]{G6}. The error in a bounded Lipschitz
test is bounded by its Lipschitz constant times the clock error,
plus twice its supremum norm times the omitted mass, plus a term
tending to zero. Exhaust the entry prefixes and let \(\varepsilon\)
and \(\eta\) tend to zero. This yields fine-sequence weak convergence
to \(\one_D(y)w(y)\,dy\,Q_\beta(y,dt)\,\delta_{\cen}(dz)\), with mass
\(m(D)\).

\section{The modulus of the joint law}\label{sec:classes}

\subsection{One modulus at all depths}\label{sec:onemod}
Before the first permanent section there are finitely many post-entry
record types. Indeed the passage to rank one has bounded length in
returns and primitive switches, and the rank-one outer graph is finite
and acyclic. Each return has finitely many guard and residue choices.
Strict leading entry prefixes of arbitrary length introduce no period:
their integral affine maps fill the same six entrance lattices.
For a homogeneous rank-one arrival, \(NT=0\) and \(T=\tau u\).
The least multiple of \(P\) restoring the full centered coset is
\[
 P_{\rm cent}=P\ell_0/\gcd(\ell_0,\tau).
\]

\begin{lemma}[Centred autonomy]\label{lem:autonomy}
\textup{(a)} On paired states the exact update, written in
\(X=6q-B\vct\), is a map \(X\mapsto\Psi(X)\) that does not depend on \(B\)
\cite[Theorem~4]{G4}; see \cite[Section~2]{G7} for the eleven-digit formula.
\textup{(b)} Consider a rank-one return certified by Lemma~\ref{lem:lines},
with homogeneous incoming line and with outgoing leading interval whose
affine hull contains zero, or with central endpoint. Write its winning guard
as \(a_GX+b_GB+c_G\) with decrement \(s>0\). Then \(b_G=0\), and a
competing guard \(H\) tied identically in leading time has
\(b_H/s_H=b_G/s=0\). Hence the exit residue \(G\bmod s\) is a function of the
centred state alone.
\end{lemma}

\begin{proof}
(a) In centered units the two digits in borrowed level \(i\) have
common base coefficient \(v_i\), and offsets
\(X_i-6\one_{i=5}\) and \(-X_{6-i}-6\one_{i<5}\). On a paired state the
singleton and the five levels occur in a fixed order, so the centered
output depends only on these offsets. This uses pairing only, not
permanence. At the outer sections, which need not be permanent, pairing
is supplied by the continuation argument in the proof of
Proposition~\ref{prop:onemod}.

(b) The leading exit satisfies \(a_G x'=-b_G\), where
\(x'=\lim X'/B\). The affine hull of the outgoing interval or its
central point contains zero, so \(b_G=0\). A competing guard \(H\)
tied identically in leading time has \(b_H/s_H=b_G/s=0\), by evaluating
the common form at zero on the homogeneous incoming line. Every section of
Lemma~\ref{lem:lines} lies on a homogeneous leading line
(Section~\ref{sec:collar}), so its certified returns satisfy the hypothesis. The intercept
comparison uses the exact normal offset.
\end{proof}

\begin{proposition}[One modulus at all depths]\label{prop:onemod}
Let \(P_*\) be a common multiple of the compatibility periods of all leaves
through first rank one, including central leaves, and of the centred
periods \(P_{\rm cent}\) of all first rank-one arrivals; there are finitely
many. Fix a class \(B=B_0+kP_*\), \(k\ge0\), and an arrival record with chart
\(X=Mq_0+B\mathbf b+c\) and admissible original inputs \(q_0\in p_B+\Lambda\)
(Section~\ref{sec:returns}). The pairs \((B,q_0)\) allowed by the record's
unrestricted congruences form an affine lattice in \(\Z\times\Z^5\). There
are a vector \((P_*,a)\) of its direction lattice, an admissible origin
\(p_0\) at \(B_0\), and \(\sigma\in\Lambda\) such that, writing the
admissible original inputs at \(B_0+kP_*\) as \(q_0=p_0+k(a-\sigma)+\mu\)
with \(\mu\in\Lambda\), the centred state at first rank-one arrival is
\(M(p_0+\mu)+B_0\mathbf b+c\), independent of \(k\). No original indices change; the original cell,
position and elapsed prefix are read at the actual \(q_0\). At every collar
depth, each later residue split pulls back to the same congruence on the
free coordinate \(\mu\) for all \(B\) in the class.
\end{proposition}

\begin{proof}
In the chart \(X=Mq_0+B\mathbf b+c\) of Section~\ref{sec:returns} at
first rank-one arrival, the original integer input satisfies
\(q_0\in p_B+\Lambda\). On an arrival record the unrestricted joint
congruences in \((B,q_0)\) form an affine lattice whose base projection
contains the progression \(B_0+kP_*\). Choose an integral lift
\((P_*,a)\) in its direction lattice. The admissible original cosets are
then \(p_0+ka+\Lambda\), where \(p_0\) is an admissible origin at \(B_0\),
and the centered displacement of these origins is
\(\Delta_a=Ma+P_*\mathbf b\). Since \(P_*\) is a multiple of the compatibility
periods and of \(P_{\rm cent}\), the centered coset \(Mp_B+B\mathbf b+c+M\Lambda\) is restored after
\(P_*\) (here \(M\Lambda\) is the inherited tangent lattice \(\ell_0u\Z\) of
Section~\ref{sec:collar}), so \(\Delta_a\in M\Lambda\); choose
\(\sigma\in\Lambda\) with \(M\sigma=\Delta_a\). Writing
\(q_0=p_0+k(a-\sigma)+\mu\) with \(\mu\in\Lambda\) gives
\(X=M(p_0+\mu)+B_0\mathbf b+c\) for every \(k\). In this free coordinate the centered state and every subsequent
autonomous congruence are independent of \(k\). The change of origin
preserves every original index, so at each fixed depth the weighted Riemann
sum in \eqref{eq:transfer} retains the original position and the prefix
clock as marks. (Holding \(q_0\) fixed instead need not give an admissible
input at the new base.)

The paired argument starts at first rank one. At the leading
exit of a long return, the endpoints of its last complete copy have
the same limit because they differ by \(d/B\). The exit is then
periodic for \(K\), by continuity. Bounded switches follow phases
of that orbit, and the next long word covers a period with paired
phases by Lemma~\ref{lem:lines}. Each adjacent interlevel slack has a
positive central value. Its zero set on a homogeneous line is
proper, and the same is true of the zero-inner guard. Thus the
first short update is strictly interlevel paired away from a proper
wall. The corresponding leading branch is homogeneous about the
centre and has unimodular derivative. It sends a nonzero interval
to a homogeneous line of nonzero variation. Induct through the
short prefix. This gives paired continuation at positive rank-one
radii, including the outer sections.

By Lemma~\ref{lem:autonomy}, the paired centered map does not depend on
\(B\), and every later winning guard has \(b_G=0\). Every later
congruence therefore pulls back on the same original free lattice, so
one modulus suffices for every collar depth.
\end{proof}

\subsection{The disappearance of the factor eleven}\label{sec:eleven}
The compatibility and centered periods at first rank one are
\((6,6),(6,12),(12,12),(6,66)\)
(Lemma~\ref{lem:periods}). Compatibility through the first rank-one
arrival, including its original residue splits, uses only six and
twelve. The centered period divides 132. On a twelve-class with
\(P_{\rm cent}\mid12\) the normalized original marks and the entire
clock suffix agree by the construction of Proposition~\ref{prop:onemod} applied to
these records with \(P_*=12\); its proof uses only the periods of the record
at hand. A change of
representative within a fixed progression changes a matched input
by a bounded translation over a fixed shift, and changes the
physical elapsed prefix by only a bounded number of updates.
The parametrization on the whole progression gives equality of
the marked limiting weights at every fixed depth.

For \(P_{\rm cent}=66\), the eleven refined classes over a twelve-class
have the same normal sheet and different tangent phases. On every
relevant strict interval the first rank-one macroscopic return is
terminal. This assertion is checked with the full physical tangent,
which contains all eleven phases. The short route and winning guard
are determined by the leading interval and, in structural ties,
the common normal offset. The return's leading duration is
\(\ell\alpha/s\) in every phase. Its final remainders move the actual
endpoint only within a bounded centered set. The section comparison
controls first capture if it occurs within that return. Summing
the final remainder probabilities restores the full incoming
weight, at the same original position and with the same total
clock. Thus the factor eleven disappears even jointly with the
original position and elapsed time.

There are 3,248 period-66 arrivals in each even root and none in the
odd roots. Their 49 normal lines per even root give 98 sections
after the split at the permanence boundary. One explicit example
has \(B=6k\) and
\[
 X=(6,18,-6,-30,-24)+\lambda(1,1,0,-1,-1),\qquad \lambda\equiv6k\pmod{66}.
\]
On compact subintervals of \(0<\lambda/B<3/14\), for sufficiently
large \(B\), 26 switches carry this line, independently of \(B\), to the
line
\[
 X'=(-102,-192,-192,-180,-90)+\lambda(1,2,2,2,1)
\]
with the same parameter \(\lambda\). A fifteen-update word then
translates \(X'\) by \(-66(1,2,2,2,1)\), that is,
\(\lambda\mapsto\lambda-66\). The largest lower threshold among its
numerical sorting guards is \(\lambda=156\), with equality allowed
because interchanging equal digits leaves the update unchanged.
All eleven phases therefore have scaled suffix duration
\(5(\lambda/B)/22\), even in this example with decrement 66.

The projection arguments give equality of the \emph{whole} measures
\(\one_D(y)w(y)\,dy\,Q_\beta(y,dt)\,\delta_{\cen}(dz)\) for all compatible
\(\beta\) having the same residue \(r\bmod12\). This is why \(Q_r\) may be
defined, as in Section~1, from the fine system of the integer
representative \(0\le r<12\); the equality determines its conditional law
up to \(w\,dy\)-null sets. Every divergent sequence in this class has a
subsequence eventually fixed modulo \(j!\) for every \(j\), by diagonal
extraction. Applying this to any sequence violating convergence for a
bounded Lipschitz test proves convergence on the entire class.

\subsection{Minimality}
Use the bounded clock reader \(\varphi(t)=\min(1,\max(0,t))\) in two local
comparisons. Let \(y_*=(\vct+x_*)/6\) be the centre of the paired cube of
Lemma~\ref{lem:local}. Then
\[
 \int\varphi\,dQ_9(y_*,\cdot)-\int\varphi\,dQ_5(y_*,\cdot)
 =\frac{314485453}{201016442880000}>0 .
\]
At the midpoint \(y_V\) of the box of \cite[Theorem~4]{G7},
\[
 \int\varphi\,dQ_5(y_V,\cdot)-\int\varphi\,dQ_{11}(y_V,\cdot)
 =\frac{248292553}{13819880448000}>0 .
\]
The neighborhoods have strict initial entry, constant residue
weights and affine clock branches. Thus each comparison persists
on a positive-weight open set. Multiply \(\varphi\) by a nonnegative
continuous spatial bump with compact support in that set. The support
lies inside \(\Pi\), so the comparison concerns \(D\) on the coupled
input. Its integral distinguishes the corresponding joint laws.

Let \(M\) be a positive integer. If \(4\nmid M\), then
\(\gcd(M,12)\mid6\), so every \(M\)-class \(a\equiv5\pmod{\gcd(M,12)}\)
contains infinitely many bases congruent to 5 and infinitely many
congruent to 11 modulo twelve, and the measures do not converge along
it. If \(3\nmid M\), then \(\gcd(M,12)\mid4\), and every \(M\)-class
\(a\equiv1\pmod{\gcd(M,12)}\) contains infinitely many bases congruent
to 9 and infinitely many congruent to 5 modulo twelve. Hence convergence
along every class modulo \(M\) forces \(12\mid M\). Conversely, if
\(12\mid M\), each class modulo \(M\) lies in one class modulo twelve,
along which the measures converge. This completes the proof of
Theorem~\ref{th:main}.

\section{Original mass of the entrance event}\label{sec:mass}

The odd-length majorization theorem of \cite{G4} gives
\(g(Ky)\prec g(y)\). Hence the sum of the \(k\) smallest gaps is
nondecreasing for every \(k\), and the largest gap is nonincreasing.
Every leading orbit converges by phases to a periodic orbit
\cite[Lemma~2]{G6}. Almost every orbit enters the invariant open ball
\(e<1/30\) by \cite[Section~3]{G6}. Its limit remains strictly inside that
ball. The periodic catalogue there has 191 cells up to cyclic
phase, as in \cite[Appendix~A]{G6}. Lemma~\ref{lem:basins} reduces
the possible destinations of the non-entrants, up to a null set, to
the centre, a family of two-cycles denoted Two, and the period-three
family in \eqref{eq:family}. Two has gap phases \((A,A,A,F,C,A)\)
and \((A,A,C,A,F,A)\), with \(A=1/6+\zeta_1\), \(F=1/6+\zeta_2\),
\(C=1-4A-F\), \(\zeta_2\ge\zeta_1\ge0\) and \(3\zeta_1+2\zeta_2\le1/6\),
restricted to the open ball.

Since \(K(y)_1=\max(y_1,1-y_5)\), the first difference, also called
the anchor, is nondecreasing. At the centre and Two it is at most
\(5/6\). Put \(\theta_1\le\cdots\le\theta_5\) for the ordered values of
\(\min(y_i,1-y_i)\), the folded lows. The formula for \(K\) gives
\[
 (K^2y)_1=\max(1-\theta_1,\theta_4+\theta_5).
\]
The necessary second-step constraint at cutoff \(U\) is
\(\theta_1\ge1-U\), \(\theta_4+\theta_5\le U\). Exact integration
(Lemma~\ref{lem:integrals}) gives probabilities
\(5124467/80621568\) at \(U=5/6\) and
\(3882730875/17179869184\) at \(U=7/8\).
All three possible non-entrant destinations have first
difference at most \(7/8\), so the latter figure alone gives
\(m(D)\ge77.3995317809\%\).

The remaining three-cycle has successive phases
\begin{equation}\label{eq:family}
 \begin{gathered}
 (7/8,3/4,5/8,\chi,1/4),\\
 (7/8,5/8,1/2,3/8,5/8-\chi),\\
 (7/8,\chi+1/4,\chi,1/4,1/8).
 \end{gathered}
\end{equation}
Here \(3/8<\chi<1/2\). Their common gap multiset is
\(\{1/8,1/8,1/8,1/4,5/8-\chi,\chi-1/4\}\).
An orbit tending to this family has all first differences at most
\(7/8\). Applying the first-difference formula to the next step
gives both end gaps at least \(1/8\) at time three. Majorization
gives the two other constraints defining \(E_3\) because the
limiting sum of the three smallest gaps is \(3/8\) and the
limiting largest gap is \(1/4\). This includes trajectories
switching between nearby sorting branches. By the union bound,
\[
 m(D^c)\le \Prob((K^2y)_1\le5/6)+m(E_3).
\]
The last exact integration of Lemma~\ref{lem:integrals} yields
the numerical bound in Proposition~\ref{prop:mass}, with
directed rounding. The exact lower bound is
\(92.9921494990216668\ldots\%\).

For comparison, direct integration through step four gives
\(13.3940217689768\ldots\%\), hence the downward bound
\(13.3940217689\%\). Disjoint certified boxes of first entry
at steps 5 through 19 add exactly
\(16984163369908523775525/604462909807314587353088\) in
original weight. This gives the earlier-entry lower bound
\(16.2038159009\%\). The stronger bound above uses the
non-entrant union directly.

\section{Scope and finite-scale evidence}\label{sec:scope}

\emph{Conditioning, observable and quantifier.} Theorem~\ref{th:main}
concerns the unnormalized measures obtained by restricting to
\(\{y\in D\}\), of limiting mass \(m(D)\); the conditional law given
entrance is obtained by dividing by this mass. The event concerns the
real digits that couple all bases. To substitute the word event
\(q_B/B\in D\) it would suffice that the two indicators disagree with
probability tending to zero; the estimate \(\|q_B/B-y\|_\infty<1/B\) gives
this if \(m(\partial D)=0\), which is not proved here.
Section~\ref{sec:kernel} works directly with the coupled event. The theorem gives no rate of convergence. The twelve
limits are not claimed to be pairwise distinct: Section~\ref{sec:classes}
separates only class 9 from class 5 and class 5 from class 11, each on a
small neighbourhood. The kernel \(Q_r\) is explicit only locally, as a
finite mixture of affine clocks (Table~\ref{tab:cube} and
\cite[Table~2]{G7}); globally it is the law of an infinite random route.
Internal labels may change under twelve-shifts without changing the
measure. For example on the box of \cite{G7} two equal-weight checkpoint
residues, 603 and 639, exchange under such a shift and give the same
pooled measure.

The common modulus in Theorem~\ref{th:main} is a property of the
\emph{joint} law on \(D\). Whether the clock marginal on \(D\)
retains that modulus remains open. In four of the twelve classes the
clock has an explicit bound, which is not used for
Theorem~\ref{th:main}.

\begin{proposition}\label{prop:fourlift}
Let \(r\in\{3,5,9,11\}\). For \(w\,dy\)-almost every \(y\in D\), the
kernel \(Q_r(y,\cdot)\) is supported in \([0,(6Y_1-5)/3]\subset[0,1/3)\),
where \(Y=K^{j(y)}y\) is the entrance point.
\end{proposition}

\begin{proof}[Proof sketch]
The certified decrease of the leading centered first coordinate per unit
scaled time is at least six through rank three and at least three at
ranks two and one. These minima, \(6,6,6,3,3\) at ranks five through one,
were checked exhaustively on the primary cover for roots 3 and 5 modulo
six, which are the four lifts (Section~\ref{sec:verification}). Bounded
switches preserve the leading first coordinate because they follow a
periodic orbit whose first coordinate is nondecreasing. The terminal
value in centered coordinates is zero. Summing from the entrance
\(Y\in\Pi\), where the first centered coordinate is \(6Y_1-5\), gives
\(0\le t\le(6Y_1-5)/3<1/3\).
\end{proof}

A prospective physical experiment near base \(3.072\cdot10^6\)
used 36,000 independent real inputs shared across its arms. It measured
the clipped clock \(\varphi(T_B/B)\) times the indicator for leading entry by
step 500 and physical capture within the run. The equally weighted
class-3 minus class-5 comparison, pooling the shifts by 12, was
\(2.54\cdot10^{-4}\), with paired pointwise 95\% interval
\([2.23,2.86]\cdot10^{-4}\). This is finite-base evidence concerning
the clock marginal. The local tests used for minimality integrate
against spatial weights and give no spatially integrated sign.

Separately, in 30,000,000 independent draws the leading-map experiment
had seven unresolved grid trajectories at step \(2^{20}\).
Allowing all 28 ambiguous draws from the grid-to-continuum comparison
gives at most 35 surviving real draws at that horizon. The
one-sided 95\% binomial upper confidence bound for the real-input
survival probability at that horizon is \(1.54681\cdot10^{-6}\).
Since \(D^c\) lies in that survival event, this gives the statistical
lower statement \(m(D)\ge1-1.55\cdot10^{-6}\) at 95\% confidence.
Proposition~\ref{prop:mass} is the deterministic bound.

The full joint law on \(D^c\) is open. For leading inputs tending to
the centre, the same finite leading shadowing and centered clock
estimate give scaled time zero and terminal mark \(\cen\). The
weight of this basin is unknown. Whether the leading basin of Two has
positive weight is not settled here; a law on that basin would need a
transfer to the physical return sheets. The measure of the switching
basin of \eqref{eq:family} is also open.

\section{Computations and verification}\label{sec:verification}

From a fresh copy of the ancillary directory, run
\texttt{bash reproduce.sh}. The script verifies the input hashes,
runs in a private tree and compares exact results with the supplied
records. The box leaf stream is regenerated and checked by SHA-256.
Arithmetic uses integers and rational fractions. The verification
environment uses Python~3.9.6, Clang~21.0.0, SymPy~1.14.0, and
Yices~2.7.0 for the singular-ray linear query. The two-solver
centered-clock checks are those of \cite[Appendix~E]{G6}. The
empirical readers in Section~\ref{sec:scope} use their frozen sample plans.

\begin{center}
\newcolumntype{L}[1]{>{\raggedright\arraybackslash}p{#1}}
\renewcommand{\arraystretch}{1.15}
\begin{tabular}{@{}L{.25\textwidth}L{.33\textwidth}L{.36\textwidth}@{}}
\toprule
Result & Main check & Separate coverage \\
\midrule
Finite routes and weights &
Exact guards, residue indices, cube and box controls &
Separate borrowed reader and source-lattice controls \\
Rank passage (Lemma~\ref{lem:ingress}) &
Six-root inherited charts &
Separate clipping, reader and lattice transport \\
Collar access (Lemma~\ref{lem:lines}) &
Six normal-sheet graphs &
Independent closure at all six roots \\
Arithmetic (Lemma~\ref{lem:periods}) &
All ingress periods and exceptional returns &
A separate lattice traversal and a separate physical audit \\
Local separations (Lemma~\ref{lem:local}) &
Exact cube clocks and full box weights &
Separate lattice reductions and local means \\
Basin exclusions (Lemma~\ref{lem:basins}) &
Exact normal and guard-cone exclusions &
Separate unpaired sorting and elimination, using earlier inputs \\
Mass integrals (Lemma~\ref{lem:integrals}) &
Exact inversion and density integration &
Separate anchor and three-step integrations \\
Four-lift clock bound (Proposition~\ref{prop:fourlift}) &
First-coordinate drift on the primary cover &
Separate graph closure \\
Physical panels &
Frozen original-word experiments &
Direct-subtraction controls \\
\bottomrule
\end{tabular}
\end{center}

\appendix
\section{Finite geometric and arithmetic certificates}\label{ap:certificates}

\begin{lemma}\label{lem:ingress}
The closure of \(\Pi\) has 77 vertices, and its only maximizer of the
energy is the gap vector \((0,1,1,1,1,1)/5\), with \(e=1/30\). Fix a root. Off a finite
union of proper walls, \(\Pi\) is partitioned into finitely many leaves,
each an open polyhedral cell \(C\) together with a coset of finite index
in the entrance lattice, obtained by the residue splits along its route,
with the following property. For each leaf and each compact
\(\mathcal K\subset C\) there is \(B_0(\mathcal K)\) such that, for
\(B\ge B_0(\mathcal K)\) on the leaf's recorded compatible base progression
in that root, every integer state on the leaf's coset (the image of
\(p_B+\Lambda\) under the entry chart) on that progression whose normalized position lies in \(\mathcal K\) follows the leaf's
certified record to first rank one: at most six macroscopic translating
returns, each with its recorded winning guard and exit residue; at most
230 primitive switches in the even roots and 246 in the odd roots; and
rank history \(54321\), \(543221\) or \(5432221\). Over the six roots the
partition has 22,110 third exits and 1,292,015 inherited first rank-one
arrivals.
\end{lemma}

For the chamber check, impose the fifteen within-level and the
adjacent interlevel rows of the five unrolled phases on the gap
simplex. Enumerate vertices by exact clipping and evaluate the convex
energy at every vertex. The physical route traversal starts with
the resulting open cell and one root coset. Its worklist retains
the uncovered side of each proper cut. A rational interior point
proposes a word. The verifier checks all of its primitive sorting
and interlevel rows, checks that the affine word fixes the full
leading chart and translates the incoming sheet, and retains only
the subcell certified by these rows. Strict rows and structural
equalities are read with the integer offsets described in
Section~\ref{sec:routes}. For each leading first winner, it clips by
all competing time rows and takes each allowed integer residue
with its exact exit map. On rank-two cells, projecting to two
independent coordinates and taking the exact convex hull leaves
the same leading cell. The inherited lattice remains unchanged
by the coordinate choice. Exhaustion of the worklist supplies
coverage. Interior witnesses, short allowances and word-length
allowances are search parameters. Coverage follows from the
complete checked partition and the absence of unfinished cells.
The uniformity in \(B\) on compact subcells is that of
Section~\ref{sec:routes}.

\begin{lemma}\label{lem:lines}
Fix a root. Every inherited first rank-one arrival of
Lemma~\ref{lem:ingress} lies on a homogeneous leading line. Enlarge each
arithmetic tangent to the full physical tangent on the same normal sheet,
enlarge each inner nonzero interval to its whole half-ray in
\(R(X)/B<1/7\), and keep outer intervals at their actual positions. Off
proper walls the resulting sections split into finitely many strict
intervals forming a finite directed graph with the following property.
For each strict interval \(J\) and compact \(\mathcal K\subset J\) there is
\(B_0(\mathcal K)\) such that, for \(B\ge B_0(\mathcal K)\) on the
section's compatible progression, every physical state on the sheet with
leading position in \(\mathcal K\) follows a bounded short segment and then
a certified paired translating return of positive scaled duration, whose
exit, for each exit residue, lies in a recorded successor section or at
the centre. Each outer subgraph is acyclic, each inner arrow ends in an
inner section or at the centre, and the minimum of inner scaled duration
divided by incoming radius is \(5/48\) in each graph.
\end{lemma}

The independent graph counts are as follows. The short bound here
concerns one rank-one transition, separate from the ingress bound.
\[
\begin{array}{c|rrrrrr}
  \text{root}&0&1&2&3&4&5\\ \hline
  \text{sections}&26409&38997&26409&38994&26403&38977\\
  \text{edges}&164139&210713&164355&210756&164103&210528
\end{array}
\]
Use exact interval cuts to certify each chronological word and each
possible leading winner. Iterate all successors until no untested
outer interval or inner half-ray remains. On the inner rays the
duration and radius are homogeneous, so their ratio is checked on
one positive point of each certified arm. Deleting all outer
vertices of indegree zero leaves zero outer vertices. The
independent odd-root rerun allows up to 250 short updates and
reaches actual maxima 172, 171 and 169 at roots 1, 3 and 5.

\begin{lemma}\label{lem:periods}
All inherited compatibility periods through first rank one are 6 or
12. The pairs of incoming progression and centered period at first rank
one are \((6,6),(6,12),(12,12),(6,66)\). Each even root has exactly
3,248 arrivals of centered period 66, all with history \(543221\), and
the odd roots have none. Their 49 normal lines per even root give 98
canonical sections after the permanence split. On each strict interval
of these sections, read on the full physical tangent, every possible
first winning guard gives a first macroscopic return whose leading
endpoint is the centre.
\end{lemma}

Use an integral primitive tangent to compute the normal period and
the exact full centered period by the divisibility tests in
Section~\ref{sec:classes}. The independent ordinary-difference
implementation checks restoration of \(6T_q-P\vct\) by the
centered tangent lattice when \(T_q\) is the ordinary base-step
translation. It tests parallelism before computing the period.
For the exceptional sections re-read all possible first winners
on the full physical sheet. In root 0 the decrements are 66 on 56
sections, 96 on 30 and 192 on 12. In root 2 the decrement is 96 on
all 98 sections. In root 4 the decrement is 24 on 70 sections and
96 on 28. The degree-one witness in Section~\ref{sec:eleven} is
checked by sorting the eleven symbolic borrowed digits on its
stated interval.

\begin{lemma}\label{lem:local}
Put \(x_*=(2938,3225,2886,1913,387)/50000\) and \(y_*=(\vct+x_*)/6\). On
the paired cube \(\|x-x_*\|_\infty<1/100000\), where \(x=X/B\) for the
centered original input \(X=6q_B-B\vct\), every input has strict initial
entry, and for \(r=9\) and \(r=5\) the kernel \(Q_r\) is the finite mixture
of affine clocks in Table~\ref{tab:cube}. At \(y_*\) the clipped clock
means are \(240752004653/23648993280000\) for class 9 and
\(74402820149/7309688832000\) for class 5. On the box of
\cite[Theorem~4]{G7} the class-5 and class-11 means at the midpoint are
\(210475357/1784592000\) and \(108644858137/921325363200\). Both centres
are strict interior points of \(\Pi\), with finite affine clock mixtures
on neighborhoods.
\end{lemma}

\begin{table}[htbp]
\centering
\renewcommand{\arraystretch}{1.1}
\(\begin{array}{rrr|ll|l}
 n_1&n_4&d&\text{class }9&\text{class }5&t(x_*)\\ \hline
 80&35&648&2333/7776&4499/15552&4343/432000\\
 100&45&816&0&2329/20736&82079/8160000\\
 140&65&1152&11963/62208&12241/62208&12883/1280000\\
 380&185&3168&6817/62208&0&21277/2112000\\
 20&15&192&0&31/1296&781/76800\\
 80&65&792&61/1728&0&26897/2640000\\
 5&5&54&55/3888&25/1296&1843/180000\\
 65&110&918&0&11/324&7933/765000\\
 20&35&288&809/3888&841/2592&3321/320000\\
 160&295&2376&1093/7776&0&16459/1584000
\end{array}\)
\caption{Clock profiles on the cube of Lemma~\ref{lem:local}:
\(t(x)=(n_1x_1+n_4(2x_2-2x_3+x_4))/d\), with their weights in classes 9
and 5. Each weight column sums to one. At \(x_*\),
\(x_1=1469/25000\) and \(2x_2-2x_3+x_4=2591/50000\).}\label{tab:cube}
\end{table}

For the cube, verify every return and terminal leaf by the affine
guard inequalities at the vertices and the residue splits on the
inherited lattice. At both interior points every time in the cited
mixture is between zero and one. Thus the evaluations are sums of
the rational weights times the affine clock values. The box means can
be recomputed from \cite[Table~2]{G7} at the midpoint. As a separate
box control, the first three pinned variables carry index 132 and
the next two congruences carry indices 36 and 6. Each marked
partner has conditional weight \(1/28512\). This check retains the
full affine time coefficients before pooling the two partners.

\section{Basin and integration certificates}

\begin{lemma}\label{lem:basins}
Among the 191 cells of the periodic catalogue in the open energy ball,
every cell other than \(\Pi\), the centre, Two and the family
\eqref{eq:family} has a null basin under the original density.
Consequently almost every original input either reaches \(\Pi\) in
finite time or tends by phases to the centre, to Two, or to
\eqref{eq:family}.
\end{lemma}

There are 30 wholly strictly interlevel-paired catalogue families,
including the three possibilities \(\Pi\), the centre and Two.
For the other 27, verify the null-basin conditions after
a late tail is contracted about equal gaps by \(1/8\).
The contraction carries strict interlevel pairing to
\(\min_i g_i>1/7\) and preserves the within-level comparison signs.
In this region the paired return
checks exclude 18 regular boundary cells and nine other cells.
For the 18 regular boundary cells, work in the frame that
removes the cyclic permutation \(L\). Each tied comparison has
a row \(\ell_i\) and adds \(v_i\max(0,-\ell_i)\) in this frame,
with \(\ell_i(v_i)=2\). In each cell one row is nonnegative on
every active correction direction. After its scheduled comparison
this row is nonnegative and thereafter nondecreasing. A
convergent tail forces it to vanish on a fixed proper hyperplane.
In three of the nine other cells, allowing a signed reference
permutation gives the same monotone-row test. The fixed ray
has ordered radii \((2/3+2f,2/3-2f,1/3,1/3,2f)\), whose distinct
levels keep their order for \(0<f<1/6\). The verified ten-step
normal expansion from \cite[Section~3]{G6} therefore applies
throughout the ray. The remaining five cells have the following
fixed-chart checks.
\begin{itemize}
\item On the edge \(\zeta_1=0\) of Two (the zero-anchor edge), the late
  normal variables \((x,y,z)\) follow \((y+(-z)_+,|z|,x-z)\). The anchor
  offset loses \((-z)_+\) each return, where \(u_+=\max(u,0)\).
  Only one initial anchor value can converge for each fixed
  normal state, so the basin is null by Fubini.
\item On the edge \(\zeta_1=\zeta_2>0\) of Two (the positive-anchor edge),
  a nonnegative four-dimensional normal state that survives has a size
  \(W\). The survivor volume for \(W\le R\) is at most
  \(2R^4\). On the eight three-return branches, each of
  determinant of absolute value one, the bounds for the
  ratio of the new size to the old size are
  \(3/5,2/3,3/8,1/2,5/13,2/3,3/8,1/2\).
  Their fourth powers sum to less than one, making surviving
  volume zero by iterating the bounds.
\item In one period-three cell the same argument uses three
  first-return branches of determinant of absolute value one,
  ratios \(2/3,3/4,1/2\) and the sum of their cubes
  \(1457/1728\). In another, three normal variables determine
  all decrements of a fourth variable; slicing in that
  fourth variable gives nullity. In the last, use gap
  offsets \(\gamma\) and normals \(x=\gamma_2-\gamma_0\),
  \(y=\gamma_4-\gamma_3\), \(z=\gamma_4+\gamma_5-\gamma_0-\gamma_1\), and
  \(\kappa=\gamma_1+\gamma_4-2\gamma_0\).
  Then \(f=x-y+z+\kappa\) and the forms
  \(f,f-x,f+(\kappa-z)/2\) must remain nonpositive because
  they are nondecreasing and tend to zero. With \(f=-1\),
  the 64 two-return sign cylinders are infeasible. By
  homogeneity a convergent tail must have \(f=0\).
\end{itemize}
Conjugation by the contraction on the paired tail and inverse
leading branches preserve nullity. An orbit converging to a
point of the open chamber reaches it in finite time.

Of the 161 other cells, 71 admit finite unrestricted phase-labelled
normal covector sets and four have earlier local normal-map
exclusions. Three of the four earlier cases use the triangle's
normal return from \cite[Section~3]{G6}. The fourth has period
three and normal coordinates \(u=y_2+y_5-1\),
\(v=y_1-y_3-y_5\), and \(z=2y_1+y_5-2\).
The phase comparisons give, with \(k=v+u_+\),
\[
 (u',v',z')=(z-u_+-k,\ |u|+2\max(k,0),\ z+2\max(-k,0)-u_+).
\]
After one return \(v\ge0\), and thereafter \(v'\ge2v\).
If \(v_3=0\), backwards substitution gives \(u=v=z=0\)
at the starting section. Thus a convergent tail must hit the
normal plane. Another 74 close by finite guarded covector sets,
nine by joint anchor and majorization inequalities, and two by
transverse noncontraction. The remaining cell is the family
\eqref{eq:family}, indexed 181 in the catalogue. The cone check
regenerates signs and orders of the folded digits on each phase.
For a phase Jacobian \(A\), branch cone \(Gx>0\), and target entry
\((u,C)\) of a normal row and cone, test whether
\(Gx>0,\ Ax\in C\) is feasible. If so, the source certificate
contains the row \(\pm uA\) with a cone containing that preimage.
The 74 guarded families use 8,130 such entries. The normal
rows annihilate the tangents of the whole limit cell. On a
convergent tail the successive values of a seed row tend to zero
and range over a finite set of row values on the first offset.
One value must be zero. The corresponding hyperplane is
independent of the chosen limit point because the row annihilates
the tangent. Taking finite preimages preserves nullity. For
the nine inequality exclusions, the necessary upper constraints
on the first difference and the sums of smallest gaps rule out
every sufficiently long guarded word. For each noncontraction
check, some return within the certified bound (six or three
returns) has transverse norm at least that of the source.
Iteration excludes convergence from nonzero transverse norm.

\begin{lemma}\label{lem:integrals}
For the original input density \(w\), the two probabilities for
\((K^2y)_1\) at cutoffs \(5/6\) and \(7/8\) are the fractions in
Section~\ref{sec:mass}. The three-step event \(E_3\) has weight
at most \(0.65165187185\%\), rounded up. The disjoint prefix
integrations through step four and on the first-entry boxes at
steps 5--19 have the weights in Section~\ref{sec:mass}.
\end{lemma}

For the anchor integrals, use the 32 original folded inverse
branches with unit absolute Jacobian. Order the five folded lows
and integrate the sum of the 32 products of six original gaps,
multiplied by \(11!\). Split at \(\theta_5=U/2\) before taking polynomial
antiderivatives under \(\theta_1\ge1-U\), \(\theta_4+\theta_5\le U\).
At \(U=1\) this gives exactly one. For the three-step event,
stratify by the unordered index triple of the three smallest
gaps and the index of the largest. This gives 60 candidate cells.
Their boundaries pull back to null sets. Invert a step to
recover the ordered lows from output differences \(q'\) by
\(\theta_i=1-q'_i-\theta_{i-1}\), \(\theta_0=0\). Clip each inverse to
the image and retain every admissible mask at each of three
steps. The affine inverse branches have unit absolute Jacobian.
There are 25 cells after image clipping, 23 of positive weight,
and 31,648 original inverse paths counted with multiplicity.

The independent integral triangulates these domains by coning
faces. For barycentric coordinates \(\xi\) on the standard
five-simplex, the integral of \(\langle z,\xi\rangle^6\) is
\((6!/11!)h_6(z)\), where \(h_6\) is the complete homogeneous
symmetric polynomial of degree six. This follows from the
monomial factorial integral on the simplex. Polarize to recover
the product of the six affine gap forms. Retain all inverse
histories with their multiplicities. The two integrations agree
exactly on every cell before the directed decimal rounding.
For the earlier entrance estimates, invert the strict chamber
through four steps and integrate its nested preimages, then
use disjoint folded boxes with all phase signs certified and
first entry between steps 5 and 19.

\section*{Data availability}
The ancillary files accompanying this paper contain the source programs,
the rational certificates, the output manifests and the replay script used
in Section~\ref{sec:verification}, with the reproduction command. They are
archived on Zenodo as version~1.8
(\href{https://doi.org/10.5281/zenodo.24396353}{doi:10.5281/zenodo.24396353}) and are available in the
\href{https://github.com/galyavin/kaprekar-dynamics/releases/tag/v1.8}{v1.8 release of the Kaprekar dynamics repository on GitHub}.

\section*{Tools}
The research, computations and drafting were carried out with substantial
assistance from AI systems (OpenAI Codex and Anthropic Claude), under the
author's direction and close supervision. The separate verification in
Section~\ref{sec:verification} was carried out in the same way, with programs
written independently of the primary ones. The computations use Python,
C, C++, SymPy and Yices. The author takes responsibility for the contents.

\begingroup
\raggedright
\begin{thebibliography}{COST}
\bibitem[COST]{COST}
E.~Chen, K.~Ono, R.~E.~Schwartz and D.~S.~Thakur,
\emph{Four-digit Kaprekar dynamics in odd bases},
J. Integer Seq. \textbf{29} (2026), Article 26.4.7.
\href{https://arxiv.org/abs/2606.20439}{arXiv:2606.20439}.

\bibitem[Dahl]{Dahl}
C.~D.~Dahl, \emph{Coarse-grained drift fields and attractor-basin entropy
in Kaprekar's routine}, Entropy \textbf{28} (2026), Article 92.
\href{https://doi.org/10.3390/e28010092}{doi:10.3390/e28010092}.

\bibitem[DZ]{DZ}
P.~Devlin and T.~Zeng, \emph{Maximum distances in the four-digit
Kaprekar process}, Integers \textbf{21} (2021), Paper A97.
\href{https://arxiv.org/abs/2010.11756}{arXiv:2010.11756}.

\bibitem[G1]{G1}
R.~Galyavin, \emph{Cycles of the four-digit Kaprekar map in even bases}, 2026.
\href{https://doi.org/10.5281/zenodo.24358739}{doi:10.5281/zenodo.24358739}.

\bibitem[G2]{G2}
R.~Galyavin, \emph{Arrival times for the four-digit Kaprekar map in even bases}, 2026.
\href{https://doi.org/10.5281/zenodo.24358744}{doi:10.5281/zenodo.24358744}.

\bibitem[G3]{G3}
R.~Galyavin, \emph{A common arrival-time law for five-digit Kaprekar dynamics}, 2026.
\href{https://doi.org/10.5281/zenodo.24358746}{doi:10.5281/zenodo.24358746}.

\bibitem[G4]{G4}
R.~Galyavin, \emph{Majorization in odd-digit Kaprekar dynamics and the nine-digit arrival law}, 2026.
\href{https://doi.org/10.5281/zenodo.24358748}{doi:10.5281/zenodo.24358748}.

\bibitem[G5]{G5}
R.~Galyavin, \emph{Seven-digit Kaprekar dynamics: global law, two clocks, parity}, 2026.
\href{https://doi.org/10.5281/zenodo.24358752}{doi:10.5281/zenodo.24358752}.

\bibitem[G6]{G6}
R.~Galyavin, \emph{Linear-time access for eleven-digit Kaprekar dynamics}, 2026.
\href{https://doi.org/10.5281/zenodo.24358476}{doi:10.5281/zenodo.24358476}.

\bibitem[G7]{G7}
R.~Galyavin, \emph{Arithmetic selection of arrival times in eleven-digit Kaprekar dynamics}, 2026.
\href{https://doi.org/10.5281/zenodo.24371697}{doi:10.5281/zenodo.24371697}.

\bibitem[KD2]{KDE}
A.~Kay and K.~Downes-Ward, \emph{Fixed points and cycles of the Kaprekar
transformation: 2. Even bases}, 2024.
\href{https://arxiv.org/abs/2408.12257}{arXiv:2408.12257}.

\bibitem[Th]{Thakur}
D.~S.~Thakur, \emph{Kaprekar phenomena}, in Topics in Number Theory,
RMS Lecture Notes Series \textbf{26}, Ramanujan Mathematical Society, 2020,
61--70. \href{https://www.sas.rochester.edu/mth/sites/dinesh-thakur/kaprekar.pdf}{Author PDF}.
\end{thebibliography}
\endgroup
\end{document}
